\documentclass[11pt]{article}
\usepackage[ngerman]{babel}
\usepackage[a4paper,margin=0.9in]{geometry}
\usepackage[T1]{fontenc}
\usepackage[utf8]{inputenc}
\usepackage{lmodern}
\usepackage{amsmath,amssymb,mathtools,mathrsfs}
\usepackage{microtype}
\usepackage{fancyhdr}
\usepackage{array,longtable,enumitem,multicol}
\usepackage{tikz}
\usetikzlibrary{patterns,calc}
\setlength{\parindent}{1.5em}
\setlength{\parskip}{0.18em}
\emergencystretch=6em
\setlength{\headheight}{14pt}
\pagestyle{fancy}
\fancyhf{}
\cfoot{\thepage}
\providecommand{\sect}[2]{\section*{\S\,#1. #2}}
\providecommand{\dd}{\,d}
\newcommand{\Z}{\mathbb Z}
\lhead{Weber Vol. I \S\S 124--130}
\rhead{Deutsche Quelle}
\begin{document}


\sect{124}{Reducirte Zahlen mit positiver Discriminante}

Wir haben nun eine ähnliche Untersuchung durchzuführen für ein positives $d$, also für reelle quadratische Irrationalzahlen. Eine solche Zahl
\[
        \omega=x+y\sqrt d
\]
soll reducirt heissen, wenn sie folgenden Bedingungen genügt:
\begin{quote}
$\omega$ ist positiv und grösser als $1$, und die mit $\omega$ conjugirte Zahl $\omega'$ ist negativ und dem absoluten Werthe nach kleiner als $1$.
\end{quote}
Oder in Zeichen:
\[
\tag{1}        0<y\sqrt d-x<1<y\sqrt d+x.
\]
Deutet man in einer Ebene $x,y$ als rechtwinklige Coordinaten, so dass die Punkte der Ebene die Bilder der Zahlen $\omega$ werden, so liegen die Bilder der reducirten Zahlen in dem Theile der Ebene, der in der Fig. 27, die der Annahme $d=2$ entspricht, schraffirt ist.

\begin{figure}[h]
\centering
\begin{tikzpicture}[scale=1.05]
  \draw[->] (-2.0,0) -- (2.6,0) node[right] {$x$};
  \draw[->] (0,-0.1) -- (0,2.45) node[above] {$y$};
  \node[below] at (-1,0) {$-1$};
  \node[below] at (1,0) {$1$};
  \node[left] at (0,0.75) {$\frac1{\sqrt d}$};
  \draw[dashed] (-1.6,0) -- (1.35,1.55);
  \draw[dashed] (-.35,0) -- (2.3,1.4);
  \draw (0.05,0.85) -- (1.8,1.7);
  \draw (-0.15,1.35) -- (1.55,2.25);
  \fill[pattern=north east lines] (0.05,0.85) -- (1.8,1.7) -- (1.55,2.25) -- (-0.15,1.35) -- cycle;
  \draw (0.05,0.85) -- (1.8,1.7) -- (1.55,2.25) -- (-0.15,1.35) -- cycle;
\end{tikzpicture}
\par\smallskip{\centering Fig. 27. Das reducirte Gebiet für $d=2$.\par}
\end{figure}

Wir nehmen hier eigentliche und uneigentliche Aequivalenz zusammen und beweisen zunächst die Sätze:

\begin{quote}
1. Jede Zahl $\omega$ ist mit einer reducirten Zahl äquivalent.
\end{quote}

Man sieht dies leicht ein, wenn man $\omega$ in einen Kettenbruch entwickelt und bis zur Schlusszahl $\omega_n$ geht:
\[
\tag{2}        \omega=(\alpha_0,\alpha_1,\alpha_2,\ldots,\alpha_{n-1},\omega_n).
\]
Dann ist $\omega$ mit $\omega_n$ äquivalent, und $\omega_n$ ist, sobald $n$ grösser als $0$ ist, positiv und grösser als $1$. Da $\omega_n$ mit $\omega$ äquivalent ist, so handelt es sich nur noch um den Nachweis, dass, wenn $n$ hinlänglich gross ist, die zu $\omega_n$ conjugirte Zahl $\omega'_n$ negativ und absolut kleiner als $1$ ist. Aus (2) folgt, wenn $P_n:Q_n$ die Näherungsbrüche sind,
\[
        \omega=\frac{P_n\omega_n+P_{n-1}}{Q_n\omega_n+Q_{n-1}},
\]
oder durch Auflösung
\[
\tag{3}        \omega_n
=-\,\frac{Q_{n-1}\omega-P_{n-1}}{Q_n\omega-P_n}
=-\,\frac{Q_{n-1}}{Q_n}
-\frac{(-1)^n}{Q_n(Q_n\omega-P_n)}.
\]
Verwandeln wir in dieser Gleichung $\sqrt d$ in $-\sqrt d$, so geht gleichzeitig $\omega$ in $\omega'$, $\omega_n$ in $\omega'_n$ über, und wir erhalten aus (3)
\[
\tag{4}        \omega'_n
=-\,\frac{Q_{n-1}\omega'-P_{n-1}}{Q_n\omega'-P_n}
=-\,\frac{Q_{n-1}}{Q_n}\,
\frac{\omega'-P_{n-1}/Q_{n-1}}{\omega'-P_n/Q_n}.
\]
Ferner ist
\[
\tag{5}        \omega'_n+1=
\frac1{Q_n}\left\{Q_n-Q_{n-1}
-\frac{(-1)^n}{Q_n\left(\omega'-P_n/Q_n\right)}\right\}.
\]
Wenn nun $n$ hinlänglich gross angenommen wird, so unterscheiden sich
\[
        \omega'-\frac{P_n}{Q_n},\qquad
        \omega'-\frac{P_{n-1}}{Q_{n-1}}
\]
beliebig wenig von der nicht verschwindenden Differenz $\omega'-\omega=-2y\sqrt d$, wonach (4) zeigt, dass $\omega'_n$ negativ ist. Aus (5) aber folgt, dass $\omega'_n+1$ positiv ist, da $Q_n-Q_{n-1}$ als positive ganze Zahl mindestens gleich $1$ ist, und
\[
        \frac1{Q_n(\omega'-P_n/Q_n)}
\]
für ein hinlänglich grosses $n$ beliebig klein wird; und dadurch ist der Beweis des Satzes 1 geführt.

Da die conjugirte Zahl von der conjugirten $\omega'$ die ursprüngliche Zahl $\omega$ ist, so folgt aus der Definition der reducirten Zahlen, dass $\omega$ und $-1:\omega'$ gleichzeitig reducirt sind.

\begin{quote}
2. Ist $\omega$ reducirt und für ein ganzzahliges $\alpha$
\[
\tag{6}        \omega=\alpha+\frac1{\omega_1},
\]
so ist $\omega_1$ dann und nur dann reducirt, wenn $\alpha$ die grösste in $\omega$ enthaltene ganze Zahl ist, wenn also
\[
\tag{7}        \alpha<\omega<\alpha+1.
\]
\end{quote}

Ist die Bedingung (7) erfüllt, so ist $\omega_1$ positiv und grösser als $1$, und
\[
\tag{8}        \omega'_1=\frac1{\omega'-\alpha}
\]
ist, da $\omega'$ negativ ist, ein negativer echter Bruch; also ist $\omega_1$ reducirt. Dagegen kann $\omega_1$ nicht reducirt sein, wenn $\alpha>\omega$ ist, da dann $\omega_1$ negativ wäre, oder wenn $\alpha+1<\omega$ wäre, da sonst $\omega_1$ kleiner als $1$ wäre.

Setzt man (8) in die Form
\[
\tag{9}        -\frac1{\omega'_1}=\alpha+\frac1{-1/\omega'},
\]
so ist, da $-1:\omega'$ und $-1:\omega'_1$ auch reducirt sind, $\alpha$ die grösste in $-1:\omega'_1$ enthaltene ganze Zahl, und $\omega'$ ist nach (9) durch $\omega'_1$, also auch $\omega$ durch $\omega_1$ völlig bestimmt.

Hieraus folgt:
\begin{quote}
3. Entwickelt man eine reducirte Zahl $\omega$ in einen Kettenbruch, so sind alle Schlusszahlen $\omega_n$ wieder reducirte Zahlen, und durch eine dieser Schlusszahlen $\omega_n$ ist sowohl die folgende $\omega_{n+1}$ als die vorangehende $\omega_{n-1}$ vollkommen bestimmt.
\end{quote}

\sect{125}{Entwickelung reeller quadratischer Irrationalzahlen in Kettenbrüche}

\begin{quote}
4. Für eine gegebene Discriminante giebt es nur eine endliche Anzahl reducirter Zahlen.
\end{quote}

Die Richtigkeit dieses Satzes sieht man leicht ein, wenn man nach §. 122 die Irrationalzahl $\omega$ in die Form setzt:
\[
\tag{1}        \omega=\frac{b+\sqrt D}{2c}
=\frac{-2a}{b-\sqrt D},\qquad D=b^2+4ac.
\]
Die conjugirte Zahl ist
\[
\tag{2}        \omega'=\frac{b-\sqrt D}{2c}
=\frac{-2a}{b+\sqrt D}.
\]
Soll $\omega$ reducirt sein, so ist
\[
\tag{3}        0<\frac{\sqrt D-b}{2c}<1<\frac{\sqrt D+b}{2c}.
\]
Es muss also $\sqrt D$ dasselbe Vorzeichen haben wie $c$; nehmen wir es positiv an, so ist
\[
\tag{4}        0<\sqrt D-b<2c<\sqrt D+b.
\]
Eine zweite Form dieser Bedingungen erhalten wir aus der zweiten Darstellung
\[
\tag{5}        0<\sqrt D-b<2a<\sqrt D+b.
\]
Es sind also $a,b,c$ positiv und $b$ kleiner als $\sqrt D$.

Bedeutet $\lambda$ die grösste in $\sqrt D$ enthaltene ganze Zahl, also
\[
        \lambda<\sqrt D<\lambda+1,
\]
so hat $b$ einen der Werthe $1,2,3,\ldots,\lambda$, jedoch mit der Beschränkung, dass bei geradem $D$ nur die geraden, bei ungeradem $D$ nur die ungeraden unter diesen Zahlen für $b$ zu nehmen sind. Es ist dann $a$ und $c$ so zu bestimmen, dass
\[
\tag{6}        \frac{D-b^2}{4}=ac,
\]
und dass
\[
\tag{7}        \frac{\lambda-b+1}{2}\le c\le \frac{\lambda+b}{2},
\]
und $a$ muss zwischen denselben Grenzen liegen. Nur einer von diesen beiden Grenzwerthen ist eine ganze Zahl, der andere kann für die untere Grenze durch die nächst grössere, für die obere Grenze durch die nächst kleinere ganze Zahl ersetzt werden.

Es ist also für $b$ nur eine endliche Zahl von Werthen zulässig; dann ist $(D-b^2):4$ nur auf eine endliche Anzahl von Arten in zwei Factoren zerlegbar, und von diesen Zerlegungen sind nur die beizubehalten, in denen beide Factoren der Bedingung (7) genügen. Ausserdem sind noch solche Combinationen wegzulassen, in denen $a,b,c$ einen gemeinsamen Theiler bekommen. Darin liegt das Mittel, um für ein gegebenes $D$ alle reducirten Zahlen wirklich zu bestimmen.

Wir wollen für das Folgende die so bestimmten Zahlen $\omega$ durch das Symbol
\[
\tag{8}        \omega=\frac{\sqrt D+b}{2c}
=\frac{2a}{\sqrt D-b}=\{a,b,c\}
\]
bezeichnen.

Aus 3. §. 124 und 4. §. 125 ergiebt sich nun der folgende wichtige Satz:
\begin{quote}
5. Die Kettenbruchentwickelung einer reducirten Zahl $\omega$ ist periodisch.
\end{quote}
D. h. wenn $\omega$ in einen Kettenbruch
\[
\tag{9}        \omega=(\alpha_0,\alpha_1,\alpha_2,\ldots)
\]
entwickelt wird, so kehren immer nach einer bestimmten Anzahl von Theilnennern dieselben Theilnenner in derselben Reihenfolge wieder. Wenn die Periode aus $\nu$ Gliedern besteht, so ist
\[
\alpha_0=\alpha_\nu=\alpha_{2\nu}=\cdots,\qquad
\alpha_1=\alpha_{\nu+1}=\alpha_{2\nu+1}=\cdots,\qquad
\alpha_{\nu-1}=\alpha_{2\nu-1}=\alpha_{3\nu-1}=\cdots.
\]
Es genügt also, den ganzen Kettenbruch durch die Periode zu bezeichnen:
\[
\tag{10}       \omega=[\alpha_0,\alpha_1,\ldots,\alpha_{\nu-1}].
\]

Dieser Satz ist offenbar gleichbedeutend damit, dass die Schlusszahl $\omega_\nu$ des Kettenbruches (9) mit $\omega$ selbst identisch ist, und darin liegt auch der Beweis der Behauptung. Denn die Schlusszahlen $\omega_\nu$ gehören als äquivalente Zahlen alle zur selben Discriminante, und daher ist nach 4. die Zahl aller möglichen Schlusszahlen nur eine endliche. In der Reihe der Schlusszahlen $\omega,\omega_1,\omega_2,\ldots$ muss daher einmal eine schon dagewesene zum zweitenmal auftreten. Ist $\omega_\nu$ die erste, die zum zweitenmal auftritt, so muss $\omega_\nu=\omega$ sein, da sonst nach 3. §. 124 auch $\omega_{\nu-1}$ zum zweitenmal auftreten würde.

Es lassen sich also die sämmtlichen, zu einer Discriminante gehörigen reducirten Zahlen in Perioden anordnen:
\[
\tag{11}       \omega,\omega_1,\omega_2,\ldots,\omega_{\nu-1},
\]
so dass, wenn
\[
        \alpha_0,\alpha_1,\alpha_2,\ldots,\alpha_{\nu-1}
\]
die grössten darin enthaltenen ganzen Zahlen sind,
\[
        \omega=[\alpha_0,\alpha_1,\ldots,\alpha_{\nu-1}],\quad
        \omega_1=[\alpha_1,\alpha_2,\ldots,\alpha_0],
\]
u. s. f. ist, womit die Kettenbruchentwickelung für jede von diesen Zahlen gegeben ist.

Ist durch eine Periode das ganze System der reducirten Zahlen noch nicht erschöpft, so bildet man eine zweite Periode u. s. f. Da durch eine Zahl $\omega$ sowohl die in der Periode vorhergehende, wie die nachfolgende völlig bestimmt ist, so enthalten zwei verschiedene Perioden niemals eine gemeinschaftliche Zahl.

Von diesen Perioden gilt nun der Satz:
\begin{quote}
6. Reducirte Zahlen aus derselben Periode sind äquivalent, aus verschiedenen Perioden sind nicht äquivalent.
\end{quote}

Der erste Theil der Behauptung ist von vornherein klar, da reducirte Zahlen derselben Periode Schlusszahlen von einander sind; der zweite Theil ist durch den Satz §. 121, VI bewiesen, dass äquivalente Zahlen, bei hinlänglich weit fortgesetzter Kettenbruchentwickelung, schliesslich dieselben Schlusszahlen bekommen. Sind also beide Kettenbrüche periodisch, so müssen sie derselben Periode angehören.

Hat man für eine gegebene Discriminante $D$ nach den am Anfang des Paragraphen gegebenen Vorschriften das vollständige System der reducirten Zahlen entwickelt, so ist es leicht, diese Zahlen in Perioden zu ordnen und die Theilnenner $\alpha$ des Kettenbruches zu finden. Es sei nämlich
\[
\tag{12}       \omega=\frac{\sqrt D+b}{2c}=\frac{2a}{\sqrt D-b}
\]
eine von diesen Zahlen und
\[
\tag{13}       \omega_1=\frac{\sqrt D+b_1}{2c_1}
=\frac{2a_1}{\sqrt D-b_1}
\]
die ihr in der Periode unmittelbar folgende Zahl, so dass
\[
\tag{14}       \omega=\alpha+\frac1{\omega_1},\qquad
        \omega_1=\frac1{\omega-\alpha}
\]
ist, wenn $\alpha$ die grösste in $\omega$ enthaltene ganze Zahl bedeutet. Nach (12) und (14) ist dann
\[
        \omega_1=\frac{2c}{\sqrt D+b-2c\alpha},
\]
und die Vergleichung mit (13) giebt
\[
\tag{15}       a_1=c,\qquad b_1=2c\alpha-b.
\]
Wenn umgekehrt zwei der zu $D$ gehörigen reducirten Zahlen $\{a,b,c\}$, $\{a_1,b_1,c_1\}$ in der durch (15) dargestellten Beziehung stehen, worin $\alpha$ irgend eine ganze Zahl ist, so folgt die zweite der ersten in der Periode unmittelbar nach. Denn aus (15) folgt (14), und daher muss $\alpha$ nach 2. die grösste in $\omega$ enthaltene Zahl sein.

Die Zahl $\omega_1$ ist also aus $\omega$ durch die Bedingungen (15) vollständig und eindeutig bestimmt.

\begin{quote}
7. Man ordnet also die Zahlen $\{a,b,c\}$ von links nach rechts in der Weise, dass die letzte Zahl $c$ der vorangehenden zugleich die erste Zahl $a$ der folgenden wird, und dass die Summe der beiden mittleren Zahlen $b+b_1$ durch $2c$ theilbar ist. Diese Anordnung ist nur auf eine Art möglich, und der Quotient $(b+b_1):2c$ ist die Zahl $\alpha$, die als Theilnenner im Kettenbruch auftritt.
\end{quote}

Betrachten wir nun irgend eine primitive ganzzahlige quadratische Gleichung
\[
\tag{16}       A+B\Omega+C\Omega^2=0,
\]
in der $A,B,C$ ganze Zahlen ohne gemeinsamen Theiler sind, die der Bedingung
\[
        B^2-4AC=D
\]
genügen, so sind die beiden Wurzeln $\Omega,\Omega'$ dieser Gleichung quadratische, zur Discriminante $D$ gehörige Irrationalzahlen. Wenn wir sie in Kettenbrüche entwickeln, so werden wir nach §. 124 endlich auf Schlusszahlen $\omega_1,\omega_2$ kommen, die zu den reducirten gehören und also in den oben besprochenen Perioden enthalten sein müssen.

Es ist noch festzustellen, ob diese Schlusszahlen $\omega_1,\omega_2$ in derselben oder in verschiedenen Perioden enthalten sind. Nach 6. sind sie dann und nur dann in derselben Periode enthalten, wenn $\Omega$ und $\Omega'$ mit einander äquivalent sind. Nennen wir solche quadratische Irrationalzahlen, die mit ihren conjugirten äquivalent sind, zweiseitige Zahlen\footnote{„Zweiseitig“ brauchen wir nach Dedekind's Vorschlag statt des Gauss'schen \emph{anceps} oder des sonst üblichen \emph{ambig}; Dirichlet-Dedekind, \emph{Vorlesungen über Zahlentheorie}, 4. Auflage, 1894, S. 139.}, so können wir sagen, dass die Kettenbruchentwickelung der Wurzeln der Gleichung (16) zu einer oder zu zwei verschiedenen Perioden führt, je nachdem die Wurzeln zweiseitig sind oder nicht. Welcher von beiden Fällen eintritt, kann man an der Periode selbst erkennen.

Wenn nämlich $\omega$ eine reducirte Zahl ist, und $\omega'$ conjugirt zu $\omega$, so ist auch $-1:\omega'$ reducirt; und wenn $\Omega$ äquivalent ist mit $\omega$, so ist $\Omega'$ äquivalent mit $\omega'$, also auch mit $-1:\omega'$ (§. 120). Ist
\[
        \omega=\{a,b,c\},
\]
so ist nach (12)
\[
        -\frac1{\omega'}=\{c,b,a\}.
\]
Nach §. 124, (9) schliesst sich also $-1:\omega'$ ebenso an $-1:\omega'_1$ an, wie $\omega_1$ an $\omega$. Ist daher die Periode von $\omega$
\[
        \omega,\omega_1,\omega_2,\ldots,\omega_{\nu-1},
\]
so ist die Periode von $-1:\omega'$
\[
        -\frac1{\omega'_{\nu-1}},\ -\frac1{\omega'_{\nu-2}},\ \ldots,\ -\frac1{\omega'_1},\ -\frac1{\omega'}.
\]
Ist die Periode der Kettenbruchentwickelung von $\omega$
\[
\tag{17}       [\alpha_0,\alpha_1,\ldots,\alpha_{\nu-1}],
\]
so ist sie für $-1:\omega'$
\[
\tag{18}       [\alpha_{\nu-1},\alpha_{\nu-2},\ldots,\alpha_0],
\]
also die umgekehrte.

Ist nun $\omega$ mit $\omega'$, also auch mit $-1:\omega'$, äquivalent, so müssen nach 6. diese beiden Perioden mit einander übereinstimmen, wenn man bei einem geeigneten Gliede beginnt, oder die Periode von $\omega$ muss umkehrbar sein. Das besagt: Es müssen sich in der Kettenbruchentwickelung für $\omega$ zwei Elemente so auswählen lassen, dass die Entwickelung gleich lautet, wenn man sie von der einen dieser Zahlen nach links oder von der anderen nach rechts fortschreitend liest. In Zeichen: es muss für irgend ein passend bestimmtes $k$ und für $i=0,1,2,\ldots,\nu-1$
\[
        \alpha_i=\alpha_{\nu+k-i}
\]
sein. Wenn die Periode umkehrbar ist, so ist auch umgekehrt nach 6. $\omega$ mit $-1:\omega'$, also auch mit $\omega'$ äquivalent. Daraus folgt also das Resultat:

\begin{quote}
8. Zweiseitige Zahlen haben in ihrer Kettenbruchentwickelung eine umkehrbare Periode, und umgekehrt sind Zahlen, die in der Kettenbruchentwickelung eine umkehrbare Periode haben, zweiseitig.
\end{quote}

Schliesslich wollen wir noch bemerken, dass durch die Eigenschaft, in einen periodischen Kettenbruch entwickelbar zu sein, die reellen quadratischen Irrationalzahlen von allen anderen reellen Zahlen unterschieden sind. Denn ist $\omega$ eine in einen Kettenbruch entwickelte Zahl, und sind $\omega_n,\omega_m$ zwei Schlusszahlen dieses Kettenbruches, so ist
\[
        \omega=
\frac{P_n\omega_n+P_{n-1}}{Q_n\omega_n+Q_{n-1}}
=\frac{P_m\omega_m+P_{m-1}}{Q_m\omega_m+Q_{m-1}}.
\]
Wenn nun der Kettenbruch periodisch ist, gleichviel, ob die Periodicität gleich von Anfang beginnt oder erst im Verlaufe der Entwickelung, so ist für zwei verschiedene Werthe von $n$ und $m$, $\omega_n=\omega_m$, und man erhält also durch Elimination von $\omega_n$ eine quadratische Gleichung für $\omega$, die man in die Form setzen kann:
\[
        \left|
        \begin{array}{cc}
        P_n-\omega Q_n & P_{n-1}-\omega Q_{n-1}\\
        P_m-\omega Q_m & P_{m-1}-\omega Q_{m-1}
        \end{array}\right|=0.
\]

\sect{126}{Beispiele}

Wir wollen einige Beispiele für die Bestimmung der reducirten Zahlen und Perioden hier durchführen, um die Anwendung der Methode zu zeigen; die Beispiele lassen sich natürlich ganz nach Belieben vermehren.

\paragraph{1. $D=29$.}
Die grösste in $\sqrt D$ enthaltene ganze Zahl $\lambda$ ist hier $5$; also kann $b$ nur die Werthe $1,3,5$ haben, woraus
\[
        ac=\frac{D-b^2}{4}=7,5,1.
\]
Die Grenzen, in denen $a$ und $c$ liegen müssen, sind nach §. 125, (7), für die drei Werthe von $b$:
\[
        3,3;\qquad 2,4;\qquad 1,5.
\]
Man erhält also nur die einzige reducirte Zahl
\[
        \{1,5,1\}.
\]
Die Periode besteht aus einem einzigen Gliede, und der Kettenbruch für $\omega$ ist
\[
        \omega=\frac{5+\sqrt{29}}2=(5,5,5,\ldots).
\]

\paragraph{2. $D=116=4\cdot29$.}
Hier ist $\lambda=10$ und $b$ hat einen der Werthe
\[
        b=2,4,6,8,10,\qquad
        ac=28,25,20,13,4.
\]
Die Grenzen für $a$ und $c$ sind
\[
        5,6;\quad 4,7;\quad 3,8;\quad 2,9;\quad 1,10.
\]
Nach Weglassung der Fälle, in denen $a,b,c$ einen gemeinschaftlichen Factor $2$ haben, ergeben sich die reducirten Zahlen
\[
\{5,4,5\},\quad \{4,6,5\},\quad \{5,6,4\},\quad
\{1,10,4\},\quad \{4,10,1\}.
\]
Sie ordnen sich in eine Periode, wobei das erste Glied am Ende noch einmal zugesetzt ist:
\[
\{4,10,1\},\ \{1,10,4\},\ \{4,6,5\},\ \{5,4,5\},\ \{5,6,4\},\ \{4,10,1\}.
\]
Für $\omega$ ergiebt sich die Kettenbruchperiode
\[
        \frac{5+\sqrt{29}}2=[10,2,1,1,2].
\]
Beginnt man mit dem vierten Gliede, also mit $[1,2,10,2,1]$, der Entwickelung von $(2+\sqrt{29})/5$, so sieht man, dass die Periode umkehrbar ist.

\paragraph{3. $D=76=4\cdot19$, $\lambda=8$.}
\[
        b=2,4,6,8,\qquad ac=18,15,10,3.
\]
Grenzen für $a,c$:
\[
        4,5;\quad 3,6;\quad 2,7;\quad 1,8.
\]
Reducirte Zahlen:
\[
\{3,4,5\},\ \{5,4,3\},\ \{2,6,5\},\ \{5,6,2\},\ \{1,8,3\},\ \{3,8,1\}.
\]
Wir erhalten eine Periode:
\[
\{3,8,1\},\ \{1,8,3\},\ \{3,4,5\},\ \{5,6,2\},\ \{2,6,5\},\ \{5,4,3\},\ \{3,8,1\}.
\]
Die Periode des Kettenbruches für $4+\sqrt{19}$ wird
\[
        [8,2,1,3,1,2].
\]
Auch dieser Kettenbruch hat eine umkehrbare Periode, da man dieselbe Reihenfolge der Zahlen erhält, mag man von der ersten $2$ nach rechts oder von der zweiten $2$ nach links lesen.

\paragraph{4. $D=37$, $\lambda=6$.}
\[
        b=1,3,5,\qquad ac=9,7,3.
\]
Grenzen für $a,c$:
\[
        3,3;\quad 2,4;\quad 1,5.
\]
Reducirte Zahlen:
\[
        \{3,1,3\},\quad \{1,5,3\},\quad \{3,5,1\}.
\]
Periode:
\[
        \{3,5,1\},\quad \{1,5,3\},\quad \{3,1,3\},\quad \{3,5,1\};
\]
Kettenbruch mit umkehrbarer Periode
\[
        \frac{5+\sqrt{37}}2=[5,1,1].
\]

\paragraph{5. $D=148=4\cdot37$, $\lambda=12$.}
\[
 b=2,4,6,8,10,12,\qquad ac=36,33,28,21,12,1.
\]
Grenzen für $a,c$:
\[
6,7;\quad 5,8;\quad 4,9;\quad 3,10;\quad 2,11;\quad 1,12.
\]
Reducirte Zahlen:
\[
\{4,6,7\},\ \{7,6,4\},\ \{3,8,7\},\ \{7,8,3\},
\quad \{3,10,4\},\ \{4,10,3\},\ \{1,12,1\}.
\]
Hier erhalten wir drei Perioden von $1$, $3$ und $3$ Gliedern:
\[
\{1,12,1\},\ \{1,12,1\},
\]
\[
\{3,10,4\},\ \{4,6,7\},\ \{7,8,3\},\ \{3,10,4\},
\]
\[
\{4,10,3\},\ \{3,8,7\},\ \{7,6,4\},\ \{4,10,3\}.
\]
Die Kettenbruchperioden sind
\[
        [12],\qquad [2,1,3],\qquad [3,1,2].
\]
Von diesen ist die erste umkehrbar; die anderen beiden gehen durch Umkehrung ineinander über.

\paragraph{6. $D=136=4\cdot34$, $\lambda=11$.}
\[
        b=2,4,6,8,10,\qquad ac=33,30,25,18,9.
\]
Grenzen für $a,c$:
\[
        5,6;\quad 4,7;\quad 3,8;\quad 2,9;\quad 1,10.
\]
Reducirte Zahlen:
\[
\{5,4,6\},\ \{6,4,5\},\ \{5,6,5\},\ \{2,8,9\},\ \{9,8,2\},
\]
\[
\{3,8,6\},\ \{6,8,3\},\ \{1,10,9\},\ \{9,10,1\},\ \{3,10,3\}.
\]
Perioden:
\[
\{5,4,6\},\{6,8,3\},\{3,10,3\},\{3,8,6\},\{6,4,5\},\{5,6,5\},\{5,4,6\},
\]
\[
\{9,10,1\},\{1,10,9\},\{9,8,2\},\{2,8,9\},\{9,10,1\}.
\]
Kettenbruchperioden:
\[
        [1,3,3,1,1,1],\qquad [10,1,4,1].
\]
Wir haben also hier zwei verschiedene Perioden, die beide umkehrbar sind.

\sect{127}{Die Pell'sche Gleichung}

Ist $\omega$ eine zur Discriminante $D$ gehörige reducirte Zahl, so erhält man durch die Kettenbruchentwickelung
\[
\tag{1}        \omega=\frac{P_n\omega_n+P_{n-1}}{Q_n\omega_n+Q_{n-1}},
\]
und es ist $\omega_n=\omega$ immer dann, wenn $n$ ein Vielfaches von der Gliederzahl der Periode $\nu$ ist. Dann aber folgt aus (1):
\[
\tag{2}        Q_n\omega^2=(P_n-Q_{n-1})\omega+P_{n-1}.
\]
Es genüge nun $\omega$ der quadratischen Gleichung §. 122, (5):
\[
\tag{3}        c\omega^2=a+b\omega,\qquad D=b^2+4ac,
\]
und aus dieser muss, da $a,b,c$ ohne gemeinsamen Theiler sind, die Gleichung (2) durch Multiplication mit einer ganzen Zahl folgen. Bezeichnen wir diese ganze Zahl mit $u$, so ist also
\[
\tag{4}        Q_n=uc,\qquad P_{n-1}=ua,\qquad P_n-Q_{n-1}=ub.
\]
Setzen wir noch
\[
\tag{5}        P_n+Q_{n-1}=t,
\]
so dass auch $t$ eine ganze positive Zahl ist, so folgt aus (4) und (5)
\[
\tag{6}        P_n=\frac{t+ub}{2},\quad Q_n=uc,\quad
        P_{n-1}=ua,\quad Q_{n-1}=\frac{t-ub}{2}.
\]
Mit Rücksicht auf
\[
        P_nQ_{n-1}-Q_nP_{n-1}=(-1)^n
\]
ergiebt sich
\[
\tag{7}        t^2-Du^2=(-1)^n4.
\]
Die Gleichung (7) heisst die Pell'sche Gleichung. Die Aufgabe ist die, für ein gegebenes $D$ alle ihre ganzzahligen Lösungen $t,u$ zu finden. Da, wenn $t,u$ der Gleichung genügen, auch $\pm t,\pm u$ ihr genügen, so können wir uns auf die Ermittelung ihrer positiven Lösungen beschränken.

Die Formeln (4), (5) geben uns eine unendliche Zahl solcher Lösungen; wir haben nur für $n$ ein beliebiges Vielfaches der Periodenzahl $\nu$ zu wählen, für $u$ den grössten gemeinschaftlichen Theiler von $Q_n,P_{n-1},P_n-Q_{n-1}$, oder auch einfach die Zahl $Q_n:c$ zu setzen und $t$ aus (5), oder auch, nachdem $u$ bestimmt ist, direct aus (7) zu ermitteln.

Ist $\nu$ ungerade und $n$ ein ungerades Vielfaches von $\nu$, so ist
\[
\tag{8}        t^2-Du^2=-4,
\]
ist aber $\nu$ gerade oder $n$ ein gerades Vielfaches von $\nu$, so ist
\[
\tag{9}        t^2-Du^2=4.
\]
Die Gleichung (9) hat die selbstverständliche Lösung $t=2$, $u=0$; in allen anderen Lösungen von (8) oder (9) sind $t$ und $u$ von Null verschieden, und wir betrachten hier nur die positiven Lösungen.

Wir wollen nun nachweisen, dass man aus einer beliebigen Gleichung von der Form (3) durch die Formeln (4), (5) alle Lösungen von (8) und (9) erhält. Nehmen wir also an, es sei $t,u$ irgend eine positive Lösung der Gleichung (8) oder (9), und $\omega=\{a,b,c\}$ eine zur Discriminante $D$ gehörige reducirte irrationale Zahl. Wir bestimmen durch die Gleichungen (6) $P_n,Q_n,P_{n-1},Q_{n-1}$, die offenbar ganze Zahlen werden, da $t$ und $ub$ entweder beide gerade oder beide ungerade sind. Aus der Gleichung (3), der $\omega$ genügt, ergiebt sich dann (2), woraus wieder auf
\[
\tag{10}       \omega=\frac{P_n\omega+P_{n-1}}{Q_n\omega+Q_{n-1}}
\]
zu schliessen ist. Aus (8) oder (9) ergiebt sich
\[
\tag{11}       P_nQ_{n-1}-Q_nP_{n-1}=\mp1,
\]
worin, wie auch in den folgenden Formeln, das obere oder das untere Zeichen gilt, je nachdem $t,u$ die Gleichung (8) oder (9) befriedigt. Nun folgt aus (6)
\[
        Q_n-Q_{n-1}=\frac{(2c+b)u-t}{2},
\]
und weil $\{a,b,c\}$ eine reducirte Zahl ist, so ist $2c+b>\sqrt D$, also
\[
        Q_n-Q_{n-1}>
\frac{u\sqrt D-t}{2}
=\frac{u^2D-t^2}{2(u\sqrt D+t)}
=\frac{\pm2}{u\sqrt D+t}.
\]
Da aber $u$ und $t$ positive ganze Zahlen sind und $D>1$, so ist $u\sqrt D+t>2$, also $Q_n-Q_{n-1}>0$ für das obere und $>-1$ für das untere Zeichen. Weil nun $Q_n-Q_{n-1}$ eine ganze Zahl sein muss, so ist hiernach
\[
        Q_{n-1}\le Q_n,
\]
wo das Gleichheitszeichen nur im Falle des unteren Zeichens, also im Falle der Gleichung (9), möglich ist.

Ferner ist, da $b<\sqrt D$ ist,
\[
        Q_{n-1}=\frac{t-ub}{2}>
\frac{t-u\sqrt D}{2}
=\frac{t^2-u^2D}{2(t+u\sqrt D)}
=\frac{\mp2}{t+u\sqrt D};
\]
also ist $Q_{n-1}$ Null oder positiv, und Null kann nur im Falle des oberen Zeichens, also im Falle der Gleichung (8), auftreten. Daraus schliessen wir
\[
\tag{12}       0\le Q_{n-1}\le Q_n,
\]
wo das Gleichheitszeichen in der unteren Grenze nur im Falle der Gleichung (8), in der oberen Grenze nur im Falle der Gleichung (9) vorkommen kann.

Wir entwickeln nun den rationalen Bruch $P_n:Q_n$ in einen endlichen Kettenbruch
\[
        \frac{P_n}{Q_n}=(\alpha_0,\alpha_1,\ldots,\alpha_{n-1}),
\]
indem wir $n$ nach §. 115, I so annehmen, dass
\[
        (-1)^n=\mp1.
\]
Bedeutet dann $P':Q'$ den vorletzten Näherungsbruch, so ist
\[
        P_nQ'-Q_nP'=\mp1,
\]
und es ist nach §. 117, II:
\[
        0\le Q'\le Q_n.
\]
Daraus folgt nach §. 118, III:
\[
        Q'=Q_{n-1},\qquad P'=P_{n-1}.
\]
Somit ist also nach (10)
\[
        (\alpha_0,\alpha_1,\ldots,\alpha_{n-1},\omega)=\omega,
\]
d. h. $[\alpha_0,\alpha_1,\ldots,\alpha_{n-1}]$ ist eine Periode des Kettenbruches für $\omega$, aus der nach den Formeln (4), (5) die Lösung $t,u$ von (7) hergeleitet wird. Damit ist bewiesen, dass wir auf diese Weise alle Lösungen der Gleichung (7) finden.

Hieraus lassen sich noch einige wichtige Folgerungen ziehen. Aus (7) ergiebt sich, dass, wenn $u$ wächst, auch $t$ wachsen muss, und dass also, wenn $u$ einen möglichst kleinen positiven Werth hat, auch $t$ möglichst klein ist. Man kann also von einer kleinsten positiven Lösung reden, die wir mit $T,U$ bezeichnen wollen. Man erhält sie, wenn man in den Formeln (4), (5) $n$ möglichst klein, also gleich der Gliederzahl $\nu$ der Periode setzt. Ist daher $\nu$ gerade, so ist nur die Gleichung (9), nicht die Gleichung (8) lösbar; ist aber $\nu$ ungerade, so ist sowohl (8) als (9) lösbar. Da dies nur von der Discriminante $D$, nicht von der besonderen Zahl $\omega$ abhängen kann, so folgt:
\begin{quote}
9. Dass bei einer Discriminante die verschiedenen Perioden entweder alle eine gerade oder alle eine ungerade Gliederzahl enthalten.
\end{quote}

Ist $D$ gerade, also durch $4$ theilbar, so muss auch $t$ gerade sein, und die Gleichung (7) lässt sich Glied für Glied durch $4$ theilen. Ist $D$ ungerade, so sind auch $t,u$ entweder beide gerade oder beide ungerade; sind sie beide ungerade, so sind ihre Quadrate, wie alle ungeraden Quadratzahlen, nach dem Modul $8$ mit $1$ congruent, also muss $D\equiv5\pmod 8$ sein. Aber es ist nicht immer möglich, wenn $D\equiv5\pmod8$ ist, die Gleichung (7) durch ungerade Zahlen zu lösen.\footnote{Vgl. Cayley, Note sur l'équation $x^2-Dy^2=\pm4$, $D\equiv5\pmod8$, Crelle's Journal, Bd. 53; \emph{Mathematical Papers}, Vol. IV. Tafeln für die Lösungen der Pell'schen Gleichung in der Form $y^2=ax^2+1$ hat Degen berechnet (Canon Pellianus, Havniae 1817). Auch in Legendre's Zahlentheorie (deutsch von Maser) findet sich eine solche Tafel.} Ist $D\equiv1\pmod8$, so müssen $t,u$ gerade sein.

Die Beispiele des §. 126 geben, auf diese Weise behandelt, die folgenden kleinsten positiven Lösungen der Pell'schen Gleichung, wobei der Factor $4$, wo es möglich ist, weggehoben ist:
\[
\begin{array}{rcl}
5^2-29\cdot1^2&=&-4,\\
70^2-29\cdot13^2&=&-1,\\
170^2-19\cdot39^2&=&1,\\
6^2-37\cdot1^2&=&-1,\\
35^2-34\cdot6^2&=&1.
\end{array}
\]

\sect{128}{Ableitung aller Lösungen der Pell'schen Gleichung aus der kleinsten positiven}

Ist $t,u$ irgend eine Lösung der Pell'schen Gleichung, so ist der Ausdruck
\[
        \frac{t^2-Du^2}{4},
\]
der den Werth $\pm1$ hat, die Norm der beiden conjugirten Zahlen
\[
        \frac{t+u\sqrt D}{2},\qquad \frac{t-u\sqrt D}{2}.
\]
Wir wollen sie die zu $\sqrt D$ gehörigen Einheiten nennen. Sind also $(t_1,u_1)$, $(t_2,u_2)$ irgend zwei positive oder negative Lösungen der Pell'schen Gleichung
\[
        t^2-Du^2=\pm4,
\]
so haben die beiden Zahlen
\[
        \Theta_1=\frac{t_1+u_1\sqrt D}{2},\qquad
        \Theta_2=\frac{t_2+u_2\sqrt D}{2}
\]
die Norm $\pm1$, und das Gleiche gilt also auch von ihrem Product
\[
        \Theta_3=\Theta_1\Theta_2=\frac{t_3+u_3\sqrt D}{2}.
\]
Darin ist
\[
        t_3=\frac{t_1t_2+u_1u_2D}{2},\qquad
        u_3=\frac{t_1u_2+t_2u_1}{2},
\]
und $t_3,u_3$ sind, wie man sofort sieht, ganze Zahlen. Denn wenn $D$ gerade ist, so sind $t_1,t_2$ gerade, und wenn $D$ ungerade ist, so sind $t_1,u_1$ entweder beide gerade oder beide ungerade, und ebenso $t_2,u_2$. Es folgt, dass das Product zweier Einheiten wieder eine Einheit ist.

Ist $\Theta$ eine beliebige Einheit, so ist $\pm\Theta^{-1}$ die zu $\Theta$ conjugirte Einheit. Es gehören unter den von $\pm1$ verschiedenen Einheiten immer vier zusammen,
\[
        \pm\Theta,\qquad \pm\Theta^{-1},
\]
unter denen eine positiv und grösser als $1$ ist. Diese entspricht den positiven Werthen von $t,u$; eine zweite ist positiv und kleiner als $1$, und die beiden anderen sind negativ. Ist $T,U$ die kleinste positive Lösung der Pell'schen Gleichung, so ist
\[
        \Theta=\frac{T+U\sqrt D}{2}
\]
die kleinste unecht gebrochene zu $\sqrt D$ gehörige positive Einheit, und wir können nachweisen, dass in der Form $\pm\Theta^n$, worin $n$ eine ganze Zahl bedeutet, alle zu $\sqrt D$ gehörigen Einheiten enthalten sind. Es genügt dazu zu zeigen, dass der Ausdruck $\Theta^n$ für ein positives $n$ alle Einheiten, die grösser als $1$ sind, liefert. Ist $\Theta_1$ irgend eine von diesen Einheiten, so wird sie zwischen zwei auf einander folgenden Potenzen von $\Theta$ liegen:
\[
        \Theta^n\le\Theta_1<\Theta^{n+1}.
\]
Mithin
\[
        1\le \Theta_1\Theta^{-n}<\Theta.
\]
Da aber $\Theta_1\Theta^{-n}$ gleichfalls eine Einheit ist, so ist, wenn nicht das Gleichheitszeichen gilt, $\Theta_1\Theta^{-n}$ grösser als $1$ und kleiner als $\Theta$, gegen die Voraussetzung, dass $\Theta$ die kleinste unecht gebrochene Einheit sei. Es muss also
\[
        \Theta_1=\Theta^n
\]
sein.

\begin{center}
\textit{Anmerkung, die Gauss'sche Theorie der quadratischen Formen betreffend.}
\end{center}

Wir fügen, um die Verbindung unserer Betrachtungen mit der sonst bekannten Gauss'schen Theorie der quadratischen Formen herzustellen, einige Bemerkungen bei, die in unserem Zusammenhange nicht gerade erforderlich sind und daher auch übergangen werden können.

Zwei quadratische Formen
\[
        (A,B,C),\qquad (-A,B,-C),
\]
nach der Bezeichnung von Gauss, sind im Gauss'schen Sinne dann und nur dann mit einander äquivalent, wenn die Periode der Kettenbruchentwickelung aus einer ungeraden Gliederzahl besteht. Nun ist eine reducirte Zahl $\omega=\{a,b,c\}$ die erste Wurzel von
\[
        (a,\tfrac12 b,-c)\quad\text{oder}\quad(2a,b,-2c)
\]
je nachdem $D$ gerade oder ungerade ist, und die zweite Wurzel von
\[
        (-a,\tfrac12 b,c)\quad\text{oder}\quad(-2a,b,2c).
\]
Sind diese beiden Formen äquivalent, bestehen also die Perioden von $\omega$ aus einer ungeraden Gliederzahl, so entspricht jeder zweiseitigen Formenclasse ein, und jedem Paar entgegengesetzter Classen zwei Systeme äquivalenter Zahlen $\omega$; es ist also die Classenzahl im Gauss'schen Sinne ebenso gross wie die Anzahl der Perioden der reducirten Zahlen $\omega$.

Wenn aber die Zahl der Periodenglieder gerade ist, so sind die Formen $(A,B,C)$, $(-A,B,-C)$ nicht äquivalent, und die Gauss'sche Classenzahl ist doppelt so gross als die Zahl der Perioden der $\omega$.

\sect{129}{Genäherte Berechnung der reellen Wurzeln einer numerischen Gleichung durch Kettenbrüche}

Auf die Theorie der Kettenbrüche gründet sich ein von Lagrange herrührendes Verfahren, um die reellen Wurzeln einer numerischen Gleichung mit beliebiger Annäherung zu berechnen, das sich jetzt mit wenig Worten auseinandersetzen lässt.

Es sei $f(x)=0$ eine reelle Gleichung, die wenigstens eine reelle Wurzel hat. Nach einer der verschiedenen Methoden des VIII. und IX. Abschnittes können wir dann immer eine ganze Zahl $\alpha_0$ finden, so dass zwischen $\alpha_0$ und $\alpha_0+1$ eine oder mehrere der Wurzeln von $f(x)$ liegen. Transformiren wir also $f(x)$ durch die Substitution
\[
        x=\alpha_0+\frac1{x_1}
\]
in $f_1(x_1)$, so wird $f_1(x_1)$ eine oder mehrere Wurzeln haben, die grösser als $1$ sind. Wir bestimmen also eine positive ganze Zahl $\alpha_1$, so dass zwischen $\alpha_1$ und $\alpha_1+1$ eine oder mehrere der Wurzeln von $f_1(x_1)$ liegen, und transformiren $f_1(x_1)$ durch die Substitution
\[
        x_1=\alpha_1+\frac1{x_2}
\]
in $f_2(x_2)$. So fährt man fort. Schliesslich wird man auf diese Weise dazu kommen, eine einzige Wurzel zu isoliren, und man erhält in der Reihe der Näherungsbrüche
\[
        (\alpha_0),\quad(\alpha_0,\alpha_1),\quad
        (\alpha_0,\alpha_1,\alpha_2),\quad\ldots
\]
eine Reihe rationaler Zahlen, die abwechselnd über und unter einer der Wurzeln $x$ liegen und sich dieser Wurzel unbegrenzt annähern. Ist man bis zu dem Näherungsbruch
\[
        (\alpha_0,\alpha_1,\ldots,\alpha_{\nu-1})=\frac{P_\nu}{Q_\nu}
\]
vorgedrungen, so unterscheidet sich dieser von dem wahren Werth von $x$ um weniger als $1:Q_\nu^2$. Die Grösse des Nenners giebt also unmittelbar ein Maass für die bereits erreichte Genauigkeit.

In den gewöhnlichen Fällen wird man die Zahlen $\alpha_0,\alpha_1,\alpha_2,\ldots$ einfach dadurch erhalten, dass man auf die Zeichenwechsel der Functionen $f,f_1,f_2,\ldots$ achtet.

Um dies Verfahren an einem Beispiel zu erproben, nehmen wir die cubische Gleichung, die wir schon früher betrachtet haben,
\[
        x^3-2x-2=0.
\]
Sie hat eine Wurzel zwischen $1$ und $2$, also haben wir $\alpha_0=1$ zu setzen. Man erhält, wenn man die oben angegebene Transformation ausführt, die folgende Kette von Gleichungen:
\[
\begin{array}{rcll}
x^3-2x-2&=&0,& \alpha_0=1,\\
3x^3-x^2-11x-1&=&0,& \alpha_1=1,\\
2x^3-4x^2-8x-3&=&0,& \alpha_2=3,\\
9x^3-22x^2-11x-2&=&0,& \alpha_3=2,\\
43x^3-6x^2-32x-9&=&0,& \alpha_4=1,\\
x^3-94x^2-132x-46&=&0,& \alpha_5=95,\\
3561x^3-9183x^2-191x-1&=&0.&
\end{array}
\]
Also
\[
        x=(1,1,3,2,1,95,\ldots)
\]
und die Näherungsbrüche sind
\[
        1,\quad \frac21,\quad \frac74,\quad \frac{16}{9},\quad
        \frac{23}{13},\quad \frac{2201}{1244},\quad \frac{4425}{2501}.
\]
Wenn wir die beiden letzten Brüche in Decimalbrüche verwandeln, so erhalten wir für $x$ die beiden Grenzen
\[
        1{,}7692926,\qquad 1{,}76929228,
\]
und der Fehler in der letzteren, etwas zu kleinen Zahl ist kleiner als
\[
        0{,}00000016.
\]
Dass man hier nach wenig Schritten ein gutes Resultat erhält, beruht auf dem Umstände, dass hier ziemlich früh ein grösserer Theilnenner $95$ auftritt. In den meisten Fällen bekommt man für die Theilnenner kleine Zahlen; dann wachsen die Nenner der Näherungsbrüche langsam, und man erhält nur mühselig ein Resultat von grösserer Genauigkeit.

Die Auffindung grösserer Theilnenner, z. B. $95$, wird durch die Bemerkung erleichtert, dass, wenn eine Wurzel einer Gleichung einen grossen Werth hat, der negative Coefficient der zweithöchsten Potenz der Unbekannten eine, wenn auch nur rohe Annäherung an diese Wurzel giebt, wie z. B. in der vorletzten der obigen Gleichungen $94$ an $95$.

\sect{130}{Rationale Wurzeln ganzzahliger Gleichungen. Reducible Gleichungen}

Wir schliessen diesen Abschnitt mit einer Betrachtung über ganzzahlige Gleichungen, die, wenn sie auch nicht unmittelbar mit der Theorie der Kettenbrüche in Beziehung steht, doch hier am passendsten eine Stelle findet.

Wenn man es mit Gleichungen zu thun hat, deren Coefficienten rationale Zahlen sind, so wird man zunächst nach etwaigen rationalen Wurzeln suchen. Wenn die Gleichung
\[
        f(x)=a_0x^n+a_1x^{n-1}+a_2x^{n-2}+\cdots+a_{n-1}x+a_n=0
\]
rationale Coefficienten hat, so können wir immer annehmen, dass diese Coefficienten ganze Zahlen sind; man hat nur nöthig, die ganze Gleichung mit einem gemeinschaftlichen Vielfachen aller Nenner zu multipliciren. Wir können aber auch noch weiter annehmen, dass $a_0=1$ sei; denn setzen wir in dem Product $a_0^{\,n-1}f(x)$
\[
        a_0x=X,
\]
so wird die Gleichung
\[
        X^n+a_1X^{n-1}+a_0a_2X^{n-2}+\cdots+a_0^{\,n-2}a_{n-1}X+a_0^{\,n-1}a_n=0.
\]
Wir wollen also jetzt annehmen, dass in
\[
\tag{1}        x^n+a_1x^{n-1}+a_2x^{n-2}+\cdots+a_{n-1}x+a_n=0
\]
die Coefficienten $a_1,a_2,\ldots,a_n$ ganze Zahlen sind. Eine rationale Wurzel von (1) kann nicht eine gebrochene Zahl sein. Denn nehmen wir an, es werde (1) befriedigt durch
\[
        x=\frac pq,
\]
wo $p$ und $q$ ganze Zahlen ohne gemeinsamen Theiler sind und $q$ positiv und grösser als $1$ ist, so folgt durch Multiplication mit $q^{n-1}$
\[
        p^n:q+a_1p^{n-1}+a_2qp^{n-2}+\cdots+a_nq^{n-1}=0.
\]
Es müsste also $p^n:q$ eine ganze Zahl sein, was unmöglich ist. Wenn nun (1) dadurch befriedigt werden kann, dass man für $x$ eine ganze Zahl setzt, so muss diese ganze Zahl nothwendig ein Theiler von $a_n$ sein, und man hat also nur die verschiedenen Divisoren von $a_n$, mit positivem und negativem Zeichen behaftet, versuchsweise einzusetzen.

Auf die Frage nach rationalen Wurzeln einer ganzzahligen Gleichung kommt auch die allgemeine Frage zurück, ob eine ganze rationale Function
\[
        f(x)=x^n+a_1x^{n-1}+a_2x^{n-2}+\cdots+a_{n-1}x+a_n
\]
durch eine rationale Function von niedrigerem Grade $\nu$,
\[
        \varphi(x)=x^\nu+\alpha_1x^{\nu-1}+\alpha_2x^{\nu-2}+\cdots+\alpha_{\nu-1}x+\alpha_\nu,
\]
mit rationalen Coefficienten theilbar sein kann. Wenn $f(x)$ durch $\varphi(x)$ theilbar ist, so ist der Quotient $\psi(x)$ wieder eine ganze rationale Function mit rationalen Coefficienten vom Grade $n-\nu$. Da $f(x)$ auch durch $\psi(x)$ theilbar ist, so können wir bei der Entscheidung unserer Frage $\nu\le \frac12n$ annehmen. Ausserdem folgt aus §. 2, dass die Coefficienten $\alpha$ und $\beta$ von $\varphi$ und $\psi$ ganze Zahlen sein müssen.

Um nun über die Möglichkeit einer solchen Theilung zu entscheiden, nehme man die Function $\varphi(x)$ mit unbestimmten Coefficienten $\alpha$, und dividire $f(x)$ durch $\varphi(x)$ nach §. 3. Es ergiebt sich dann ein Rest vom Grade $\nu-1$; setzt man dessen Coefficienten, die sämmtlich rationale Functionen der $\alpha$ sind, gleich Null, so erhält man $\nu$ Gleichungen, denen diese Coefficienten genügen müssen. Durch Elimination kann man eine Gleichung herstellen, die nur noch eine dieser Unbekannten enthält, und die dann eine ganzzahlige Wurzel haben muss. Oft ist es zweckmässiger, diese Elimination nicht wirklich auszuführen, sondern direct zu versuchen, dem System der $\nu$ Gleichungen durch ganzzahlige Werthe der $\alpha$ zu genügen.

Wir entnehmen als Beispiel eine Gleichung zehnten Grades aus der Theorie der elliptischen Functionen\footnote{Vgl. des Verfassers Werk \emph{Elliptische Functionen und algebraische Zahlen}, Braunschweig 1891. Das hier behandelte Beispiel aus der Theorie der Transformation $47$ten Grades ist der Abhandlung „Ein Beitrag zur Transformationstheorie der elliptischen Functionen“ von H. Weber (1893), Mathematische Annalen, Bd. 43, entnommen.}, bei der die Zerlegbarkeit in Factoren bestimmten Grades theoretisch feststeht:
\[
        x^{10}-36x^8+528x^6-3897x^4+14354x^2-21025=0.
\]
Wir fragen, ob die linke Seite in zwei Factoren fünften Grades zerlegbar ist. Wenn der eine der beiden Factoren
\[
        x^5+\alpha x^4+\beta x^3+\gamma x^2+\delta x+\varepsilon
\]
ist, so ist der andere, da in der gegebenen Gleichung keine ungeraden Potenzen vorkommen,
\[
        x^5-\alpha x^4+\beta x^3-\gamma x^2+\delta x-\varepsilon.
\]
Es muss also sein
\[
\begin{aligned}
x^{10}-36x^8+528x^6-3897x^4+14354x^2-21025
&=(x^5+\beta x^3+\delta x)^2\\
&\quad -(\alpha x^4+\gamma x^2+\varepsilon)^2 .
\end{aligned}
\]
Setzt man entsprechende Coefficienten einander gleich, so folgt
\[
\begin{array}{rcl}
\alpha^2-2\beta&=&36,\\
\beta^2+2\delta-2\alpha\gamma&=&528,\\
\gamma^2+2\alpha\varepsilon-2\beta\delta&=&3897,\\
\delta^2-2\gamma\varepsilon&=&14354,\\
\varepsilon^2&=&21025.
\end{array}
\]
Es sind nun ganzzahlige Werthe von $\alpha,\beta,\gamma,\delta,\varepsilon$ zu suchen. Aus der letzten Gleichung erhält man
\[
        \varepsilon=\sqrt{21025}=145=5\cdot29,
\]
und man kann $\varepsilon$ positiv annehmen. Die vierte Gleichung ergiebt, rechts nach dem Modul $145$ genommen,
\[
\tag{4}        \delta^2\equiv -1\pmod{145},
\]
und diese Congruenz hat die vier Lösungen
\[
        \delta=\pm12,\ \pm17\pmod{145}.
\]
Da $\delta$ nach der vierten Gleichung ausserdem gerade sein muss, so könnten vorläufig die Werthe
\[
        \delta=\pm12,\ \pm128,\ \pm162
\]
auftreten. Diesen Werthen entsprechend erhält man aus der vierten Gleichung für $\gamma$
\[
        \gamma=-49,\ 7,\ 41.
\]
Wenden wir noch den Modul $5$ an, so folgt
\[
        \gamma\equiv 1,-1,1\pmod5,
\]
aus der dritten Gleichung
\[
        \beta\equiv \pm2,\ \pm1\pmod5,
\]
und aus der ersten Gleichung
\[
        \alpha^2\equiv1\pm2,\ 1+1,\ 1\pm2\pmod5.
\]
Da nur $0,1,4$, nicht aber $2,3$ nach dem Modul $5$ mit einem Quadrat congruent sein können, müssen in $\delta$ die Zeichen so genommen werden:
\[
        \delta=-12,\ +128,\ -162,
        \qquad
        \alpha\equiv\pm2,\ 0,\ \pm2\pmod5.
\]
Die zweite Gleichung zeigt noch, dass der mittlere Fall auszuschliessen ist, da die rechte Seite nicht durch $5$ theilbar ist. Auch $\delta=-12$ muss ausgeschlossen werden, denn dann gäbe die zweite Gleichung mit $\gamma=-49$ die unlösbare Congruenz
\[
        \beta^2\equiv6\pmod7.
\]
Es bleibt also nur
\[
        \delta=-162,\qquad \gamma=41.
\]
Aus der dritten Gleichung folgt für $\beta$
\[
        17\beta\equiv93\pmod{145},
\]
oder
\[
        \beta\equiv -1\pmod5,\qquad \beta\equiv -1\pmod{29},
\]
also
\[
        \beta\equiv14\pmod{145}.
\]
Nimmt man $\beta=14$ an, so folgt $\alpha=-8$, und alle Gleichungen sind erfüllt. Es hat sich also die Zerlegung ergeben:
\[
\begin{aligned}
&x^{10}-36x^8+528x^6-3897x^4+14354x^2-21025\\
&\quad=(x^5-8x^4+14x^3+41x^2-162x+145)
       (x^5+8x^4+14x^3-41x^2-162x-145).
\end{aligned}
\]


\end{document}
